\chapter{Продукты, QC и воспроизводимость}

\section{Основные продукты}

Пользовательский результат обработки серии включает БД и файловые artifacts. К наиболее важным относятся:

\begin{itemize}
  \item science calibration generation;
  \item science measurements;
  \item field/object QC;
  \item candidate light curves и summaries;
  \item manifest;
  \item pipeline run log и step results;
  \item зарегистрированные immutable artifacts, если запуск выполнялся recorded workflow.
\end{itemize}

\section{Random и systematic errors}

В современной схеме случайная и систематическая составляющие ошибки разделяются. Диагностические residual floors могут присутствовать в provenance, не обязательно будучи применёнными к science error. При интерпретации следует различать measured diagnostic quantity и фактически применённую error component.

\section{Object QC}

Object QC оценивает scatter, magnitude-dependent noise, crowding/blending evidence и detection support. Эти признаки используются как характеристика качества объекта и как input для candidate logic.

\section{Field QC}

Field-level report суммирует качество серии: coverage, scatter, astrometry, aperture diagnostics, candidate statistics и другие series-level показатели.

\section{Воспроизводимость}

Для научно значимого результата следует сохранять или иметь возможность восстановить:

\begin{itemize}
  \item pipeline run ID;
  \item software version;
  \item calibration generation IDs;
  \item входную серию и identity catalog;
  \item manifest;
  \item immutable artifacts/checksums, если они опубликованы.
\end{itemize}
